% !PRINTED 194
\weberpage{194}
Das Verhalten von $H$, $D$, $Q$ bei linearer Transformation ergiebt sich aus \S.\,60, (2), (3); sind $H'$, $D'$, $Q'$ die entsprechenden Bildungen für eine transformirte Form, so hat man
\weq{8}{H' = r^2\,H,\quad D' = r^6\,D,\quad Q' = r^3\,Q.}

Die Covarianten können wir benutzen, um die cubische Form auf eine Normalform zu transformiren, die zugleich die Lösung der cubischen Gleichung giebt.

Wir wählen die Normalform
\weq{9}{f(x,y) = F(\xi,\eta) = \xi^3 + \eta^3,}
worin $\xi$, $\eta$ lineare Functionen von $x$, $y$ sind; eine Form, die, wie sich gleich ergeben wird, immer hergestellt werden kann, wenn $D$ nicht verschwindet, also die Gleichung $f = 0$ nicht zwei gleiche Wurzeln hat. Die Lösungen von $f = 0$ ergeben sich dann aus den linearen Gleichungen
\[
\xi + \eta = 0,\quad \xi + \varepsilon\,\eta = 0,\quad \xi + \varepsilon^2\,\eta = 0,
\]
worin $\varepsilon$ eine dritte Einheitswurzel ist.

Um für die Normalform (9) die Formen $H'$, $D'$, $Q'$ zu bilden, haben wir $a_0' = a_3' = 1$, $a_1' = a_2' = 0$ zu setzen und erhalten
\[
H' = 9\,\xi\,\eta,\quad D' = -\,27,\quad Q' = 27\,(\xi^3 - \eta^3).
\]

Wenn die Normalform (9) durch lineare Transformation aus der allgemeinen Form $f(x,y)$ abgeleitet ist, so ergiebt sich nach (8)
\weq{10}{
\begin{aligned}
9\,\xi\,\eta &= r^2\,H \\
-\,27 &= r^6\,D \\
27\,(\xi^3 - \eta^3) &= r^3\,Q \\
\xi^3 + \eta^3 &= f.
\end{aligned}
}
Daraus folgt
\[
r^6 = \frac{-\,27}{D},\quad 2\,\xi^3 = \frac{r^3\,Q}{27} + f,\quad 2\,\eta^3 = -\,\frac{r^3\,Q}{27} + f,
\]
also, wenn für $r^3$ aus der ersten Gleichung der Werth $-\,9 : \sqrt{-3\,D}$ gesetzt wird,
\weq{11}{
\begin{aligned}
6\,\xi^3\,\sqrt{-3\,D} &= +\,Q + 3\,\sqrt{-3\,D}\,f \\
6\,\eta^3\,\sqrt{-3\,D} &= -\,Q + 3\,\sqrt{-3\,D}\,f,
\end{aligned}
}
woraus $\xi$ und $\eta$ durch zwei Cubikwurzeln bestimmt sind, von denen nach der ersten Gleichung (10) die eine durch die andere bestimmt ist.
